
\begin{corollaire}
\label{V.1.6}
Para que $X$ sea afín, \resp un esquema,
\ifthenelse{\boolean{orig}}{es n.~y~s.}{es necesario y suficiente}
que~$Y$ lo sea.
\end{corollaire}

\begin{definition}
\label{V.1.7}
Sea
\marginpar{109}
$X$ un preesquema sobre el que un grupo finito~$G$ actúa por la derecha. Se
dice que la acción de~$G$ es \emph{admisible} si existe un morfismo
$p\colon X\to Y$ que tiene las propiedades de~\Ref{V.1.3} (lo cual implica
que $X/G$ existe y es isomorfo a~$Y$).
\end{definition}

\begin{proposition}
\label{V.1.8}
Sea $X$ un preesquema sobre el que el grupo finito~$G$ actúa por la derecha.
Para que la acción de~$G$ sea admisible es necesario y suficiente que~$X$
sea unión de abiertos afines invariantes por~$G$, o, equivalentemente, que
toda órbita de~$G$ en~$X$ esté contenida en un abierto afín.
\end{proposition}

La primera condición implica evidentemente la segunda, y esta, a su vez,
implica la primera. En efecto, sean~$T$ una órbita de~$G$ y~$U$ un abierto
afín que la contiene. La intersección de los transformados de~$U$ por
\ifthenelse{\boolean{orig}}{los~$g$ $G$}{los~$g$ de~$G$}
es entonces un abierto~$U'$ estable por~$G$, que contiene a~$T$ y está
contenido en el abierto afín~$U$. Como, dentro de~$U$, toda parte finita
posee un sistema fundamental de entornos abiertos afines, existe un entorno
abierto afín~$V$ de~$T$ contenido en~$U'$. Sus transformados por
\ifthenelse{\boolean{orig}}{los~$g$ $G$}{los~$g$ de~$G$}
son, pues, afines y están contenidos en~$U'$, que es \emph{separado}; por
tanto, su intersección~$U''$ es un abierto afín, invariante por~$G$, que
contiene a~$T$.
--- Sentado esto, la condición considerada en~\Ref{V.1.8} es
\emph{necesaria}, pues se toman las imágenes inversas~$X_i$ de abiertos
afines~$Y_i$ que recubren~$Y$. Es suficiente porque, por~\Ref{V.1.1}, pueden
construirse entonces los cocientes $Y_i=X_i/G$; en cada~$Y_i$, la imagen de
$X_i\cap X_j$ es un abierto~$Y_{ij}$ que se identifica con~$X_{ij}/G$
por~\Ref{V.1.4}. En particular, se deducen isomorfismos
$Y_{ij}\isomto Y_{ji}$ que permiten pegar los~$Y_i$ para construir~$Y$.
Serre prefiere construir directamente el espacio topológico cociente~$Y$
de~$X$ por~$G$, dotarlo del haz $p_*(\cal{O}_X)^G$ y comprobar que~$Y$ se
convierte en un preesquema y que se está entonces bajo las condiciones
de~\Ref{V.1.3}.

\setcounter{proposition}{6}

\begin{corollaire}
\label{cor:V.1.7}
Si la acción de~$G$ sobre~$X$ es admisible, también lo es la de todo subgrupo
$H$ de~$G$ (por tanto, $X/H$ existe).
\end{corollaire}

Esto puede comprobarse también directamente en la situación de~\Ref{V.1.3},
observando que siempre puede suponerse que~$X$ es afín sobre algún~$Z$ y que
los $s\in G$ actúan por $Z$-automorfismos (se toma, por ejemplo, $Z=Y$). En
efecto:

\begin{corollaire}
\label{cor:V.1.8}
Supongamos que~$X$ es afín sobre~$Z$ y que las operaciones de~$G$ son
$Z$-automorfismos. Entonces la acción de~$G$ sobre~$X$ es admisible. Si~$X$
está definido por un haz cuasicoherente~$\cal{A}$ de álgebras, $Y$ está
definido por el haz~$\cal{A}^G$ de invariantes de~$\cal{A}$ por~$G$.
\end{corollaire}

\begin{proposition}
\label{V.1.9}
Supongamos que la acción de~$G$ sobre~$X$ es admisible y que $X/G=Y$ es un
preesquema sobre~$Z$. Consideremos un morfismo de cambio de base $Z'\to Z$ y
pongamos $X'=X\times_Z Z'$, $Y'=Y\times_Z Z'$, de modo que~$G$ sigue
actuando, por transporte de estructura, sobre~$X'$, y el morfismo
$p'\colon X'\to Y'$ es invariante. Si~$Z'$ es \emph{plano} sobre~$Z$,
entonces~$p'$ sigue satisfaciendo las hipótesis
\ifthenelse{\boolean{orig}}
{de~\Ref{V.1.3}}
{de~\Ref{V.1.3},}
\ie $\cal{O}_{Y'}\to p_*'(\cal{O}_{X'})^G$ es un isomorfismo (pues~$p'$ es
afín en cualquier caso). Por consiguiente, la acción de~$G$ sobre~$X'$ es
admisible y
$(X/G)\times_Z Z'\thickapprox (X\times_Z Z')/G$.
\end{proposition}

Puede suponerse evidentemente que $Z=Y$, y la cuestión se reduce al caso en
que, además, $Y$ e~$Y'$ son afines. Hay que demostrar que, si~$B$ es el
subanillo de invariantes de la acción de~$G$ sobre~$A$ y si~$B'$ es un
$B$-álgebra plana sobre~$B$, entonces~$B'$ es la subálgebra de invariantes de
$A'=A\otimes_B B'$. Esto es inmediato, pues la sucesión exacta
\ifthenelse{\boolean{orig}}{}
{\enlargethispage{.5cm}}%
$$
0 \to B \lto{i} A \lto{j} A^{(G)}
$$
(donde el último término denota una potencia de~$A$ y $j(x)$ es el sistema
de los $s\cdot x-x$, $s\in G$) sigue siendo exacta al tensorizar por
\ifthenelse{\boolean{orig}}
{$A'$}
{$B'$}.

Debe advertirse que la hipótesis de planitud era esencial para la validez
del resultado; en particular, si~$Y'$ es un subpreesquema cerrado de
\ifthenelse{\boolean{orig}}
{$X$}
{$Y$}
(por ejemplo, incluso un punto cerrado de
\ifthenelse{\boolean{orig}}
{$X$}
{$Y$})
y~$X'$ es su imagen inversa en~$X$, entonces, en general, $Y'$ \emph{no se
identifica} con~$X'/G$. Veremos que, no obstante, sí se identifica cuando~$X$
es étale sobre~$Y$.

Para terminar, demos un formalismo tan cómodo como trivial. Sea~$Y$ un
preesquema. Como en la categoría de los preesquemas existen las sumas
disjuntas, para todo conjunto~$E$ puede considerarse el preesquema suma de
una familia $(Y_i)_{i\in E}$ de preesquemas, todos idénticos a~$Y$; este
preesquema se denotará por~$Y\times E$.
\marginpar{110}
Está caracterizado por la fórmula
\begin{equation*}
\label{eq:V.1.*}
\tag{$*$} \Hom(Y\times E,Z) = \Hom(E,\Hom(Y,Z)),
\end{equation*}
donde el segundo~$\Hom$ denota, evidentemente, el conjunto de aplicaciones
del conjunto~$E$ en el conjunto~$\Hom(Y,Z)$. Hay un morfismo canónico
$$
Y\times E\to Y
$$
que hace de~$Y\times E$ un preesquema sobre~$Y$. Como los productos fibrados
conmutan con las sumas disjuntas (en la categoría de los preesquemas), si~$Y$
es un preesquema sobre otro~$Z$, se tiene, para un cambio de base $Z'\to Z$:
$$
(Y\times E)\times_Z Z'=(Y\times_Z Z')\times E
$$
(fórmula especialmente útil si $Z=Y$). Por otra parte, se deduce
trivialmente de la definición que
$$
(Y\times E)\times F=Y\times(E\times F)
=(Y\times E)\times_Y(Y\times F)
$$
(aunque la última fórmula resulta también de la conmutatividad señalada más
arriba).

Para~$Y$ fijo, puede considerarse $Y\times E$ como un funtor de~$E$, con
valores en los preesquemas sobre~$Y$; según la fórmula anterior, este funtor
conmuta con los productos finitos. Esto permite, por ejemplo, asociar a todo
grupo ordinario~$G$ un esquema en grupos $Y\times G$ sobre~$Y$, que será
finito sobre~$Y$ si
\ifthenelse{\boolean{orig}}{$Y$ lo es}{$G$ es finito},
\ifthenelse{\boolean{orig}}{etc...}{etc.}
Más generalmente, este funtor es «exacto por la izquierda», pero no tendremos
que utilizar aquí este hecho. También conmuta trivialmente con las sumas
disjuntas y es «exacto por la derecha», como se ve de inmediato en la fórmula
de definición~\eqref{eq:V.1.*}. En particular, si el grupo finito~$G$ actúa
por la derecha sobre el conjunto~$E$, entonces actúa por la derecha sobre
$Y\times E$, y se tiene
$$
(Y\times E)/G = Y\times (E/G),
$$
donde, de hecho, el cociente del primer miembro satisface las condiciones
de~\Ref{V.1.3} (esto es inmediato).
